\subsection{BIORDERS IN THE ENVELOPING ALGEBRA OF A SPLIT REDUCTIVE LIE ALGEBRA}

Let \(\mathfrak{g}\) be a reductive Lie algebra over \(\mathbf{Q}\) , \(\mathfrak{h}\) a splitting Cartan subalgebra of \(\mathfrak{g}\) , and \(R = R(\mathfrak{g},\mathfrak{h})\) (§2, no. 1, Remark 5).

DEFINITION 1. A lattice \(\mathcal{H}\) in \(\mathfrak{h}\) is said to be permissible (relative to \(\mathfrak{g}\) ) if, for all \(\alpha \in \mathrm{R}\) , we have \(H_{\alpha} \in \mathcal{H}\) and \(\alpha(\mathcal{H}) \subset \mathbf{Z}\) .

Remarks. 1) Let B be a basis of R. A lattice \(\mathcal{H}\) in \(\mathfrak{h}\) is permissible if and only if \(H_{\alpha} \in \mathcal{H}\) and \(\alpha(\mathcal{H}) \subset \mathbf{Z}\) for all \(\alpha \in B\) .

\begin{enumerate}
\def\labelenumi{\arabic{enumi})}
\setcounter{enumi}{1}
\item
  Let \(\mathfrak{c}\) be the centre of \(\mathfrak{g}\) . Then, a lattice \(\mathcal{H}\) in \(\mathfrak{h}\) is permissible if and only if \(\mathrm{Q}(\mathrm{R}^{\vee}) \subset \mathcal{H} \subset \mathrm{P}(\mathrm{R}^{\vee}) \oplus \mathfrak{c}\) . The lattice \(\mathcal{H} \cap \mathcal{D}\mathfrak{g}\) is then permissible in the Cartan subalgebra \(\mathfrak{h} \cap \mathcal{D}\mathfrak{g}\) of \(\mathcal{D}\mathfrak{g}\) . There may exist permissible lattices \(\mathcal{H}\) such that \(\mathcal{H} \neq (\mathcal{H} \cap \mathcal{D}\mathfrak{g}) \oplus (\mathcal{H} \cap \mathfrak{c})\) (cf.~§13, no. 1.IX).
\item
  If \(\mathfrak{g}\) is semi-simple, the permissible lattices in \(\mathfrak{h}\) are the subgroups \(\mathcal{H}\) of \(\mathfrak{h}\) such that \(Q(R^{\vee})\subset \mathcal{H}\subset P(R^{\vee})\) .
\end{enumerate}

In the remainder of this number, we assume fixed a split reductive Lie algebra \((\mathfrak{g},\mathfrak{h})\) , a basis B of \(\mathrm{R}=\mathrm{R}(\mathfrak{g},\mathfrak{h})\) and, for each \(\alpha\in B\) , a pair \((x_{\alpha},y_{\alpha})\) with

\[
y _ {\alpha} \in \mathfrak {g} ^ {- \alpha}, x _ {\alpha} \in \mathfrak {g} ^ {\alpha}, [ y _ {\alpha}, x _ {\alpha} ] = H _ {\alpha}.\tag{20}
\]

If we denote by \(n_{+}\) (resp. \(n_{-}\) ) the subalgebra of g generated by the \(x_{\alpha}\) (resp. the \(y_{\alpha}\) ), we know (§3, no. 3, Prop. 9 (iii)) that

\[
\mathfrak {g} = \mathfrak {n} _ {-} \oplus \mathfrak {h} \oplus \mathfrak {n} _ {+}\tag{21}
\]

\[
\mathrm{U} (\mathfrak {g}) = \mathrm{U} (\mathfrak {n} _ {-}) \otimes_ {\mathbf {Q}} \mathrm{U} (\mathfrak {h}) \otimes_ {\mathbf {Q}} \mathrm{U} (\mathfrak {n} _ {+})\tag{22}
\]

(where \(\mathrm{U}(\mathfrak{g}),\ldots\) are the enveloping algebras of \(\mathfrak{g},\ldots\) ).

Denote by \(\mathcal{U}_{+}\) the \(\mathbf{Z}\)-subalgebra of \(\mathrm{U}(\mathfrak{n}_{+})\) generated by the \(x_{\alpha}^{(n)}\) for \(\alpha \in B\) and \(n \in \mathbf{N}\) . Let W be the Weyl group of R, \(R_{+}\) the set of positive roots relative to B.

Lemma 5. (i) \(\mathcal{U}_{+}\) is a lattice in the vector space \(\mathrm{U}(\mathfrak{n}_{+})\) .

\begin{enumerate}
\def\labelenumi{(\roman{enumi})}
\setcounter{enumi}{1}
\tightlist
\item
  For all \(\alpha \in \mathrm{B}\) , we have \(\mathcal{U}_{+} \cap \mathrm{U}(\mathfrak{g}^{\alpha}) = \bigoplus_{n \in \mathbf{N}} \mathbf{Z}x_{\alpha}^{(n)}\) .
\end{enumerate}

By definition, \(\mathcal{U}_{+}\) is generated as a \(\mathbf{Z}\) -module by the elements

\[
x _ {\varphi} ^ {(\mathbf {n})} = \prod_ {1 \leq i \leq r} x _ {\varphi (i)} ^ {(n (i))}
\]

where \(r \in \mathbf{N}\) , \(\varphi = (\varphi(i)) \in \mathrm{B}^{r}\) , and \(\mathbf{n} = (n(i)) \in \mathbf{N}^{r}\) . Give the algebra \(\mathrm{U}(\mathfrak{n}_{+})\) the graduation of type \(\mathrm{Q}(\mathrm{R})\) for which each \(\mathfrak{g}^{\alpha} (\alpha \in \mathrm{R}_{+})\) is homogeneous of degree \(\alpha\) . A monomial \(x_{\varphi}^{(\mathbf{n})}\) of the preceding type is homogeneous of degree

\[
\sum_ {1 \leq i \leq r} n (i) \varphi (i) \in \mathrm{Q} (\mathrm{R}).
\]

The monomials of this kind having a given degree q are finite in number, and generate over Q the homogeneous component of \(\mathrm{U}(\mathfrak{n}_{+})\) of degree q. This proves (i).

If \(\alpha \in \mathrm{B}\) , \(\mathcal{U}_{+} \cap \mathrm{U}(\mathfrak{g}^{\alpha})\) is contained in the sum of the homogeneous components of degrees which are multiples of \(\alpha\) ; thus, by the preceding, \(\mathcal{U}_{+} \cap \mathrm{U}(\mathfrak{g}^{\alpha})\) is generated by the \(x_{\varphi}^{(\mathbf{n})}\) such that \(\sum n(i)\varphi(i) \in \mathbf{N}\alpha\) , which forces \(\varphi(i) = \alpha\) for all \(i\) (since B is a basis of R), so

\[
x _ {\varphi} ^ {(\mathbf {n})} = x _ {\alpha} ^ {(n (1))} \dots x _ {\alpha} ^ {(n (r))} = \frac {(n (1) + \cdots + n (r)) !}{n (1) ! \dots n (r) !} x _ {\alpha} ^ {(n (1) + \dots + n (r))}.
\]

Thus, \(\mathcal{U}_{+} \cap \mathcal{U}(\mathfrak{g}^{\alpha}) \subset \bigoplus_{n} \mathbf{Z}x_{\alpha}^{(n)}\) , hence (ii).

Q.E.D.

In the remainder of this paragraph, if E and F are Z-submodules of \(U(\mathfrak{g})\) , we denote by E.F the Z-submodule of \(U(\mathfrak{g})\) generated by the products ab, where \(a \in E, b \in F\) .

PROPOSITION 3. Let \(\mathcal{H}\) be a permissible lattice in \(\mathfrak{h}\) . Let \(\mathcal{U}_{+},\mathcal{U}_{-},\mathcal{U}_{0}\) be the \(\mathbf{Z}\) -subalgebras of \(\mathrm{U}(\mathfrak{g})\) generated respectively by the elements \(x_{\alpha}^{(n)}\) ( \(\alpha \in \mathrm{B}, n \in \mathbf{N}\) ), \(y_{\alpha}^{(n)}\) ( \(\alpha \in \mathrm{B}, n \in \mathbf{N}\) ), \(\binom{h}{n}\) ( \(h \in \mathcal{H}, n \in \mathbf{N}\) ). Let \(\mathcal{U}\) be the \(\mathbf{Z}\) -subalgebra of \(\mathrm{U}(\mathfrak{g})\) generated by \(\mathcal{U}_{+},\mathcal{U}_{-},\mathcal{U}_{0}\) .

\begin{enumerate}
\def\labelenumi{(\roman{enumi})}
\item
  \(\mathcal{U}\) is a biorder in the bigebra \(\mathrm{U}(\mathfrak{g})\) .
\item
  We have \(\mathcal{U} = \mathcal{U}_{-}. \mathcal{U}_{0}. \mathcal{U}_{+}\) , \(\mathcal{U} \cap \mathfrak{h} = \mathcal{H}\) and, for all \(\alpha \in B\) ,
\end{enumerate}

\[
\mathcal {U} \cap \mathfrak {g} ^ {\alpha} = \mathbf {Z} x _ {\alpha}, \quad \mathcal {U} \cap \mathfrak {g} ^ {- \alpha} = \mathbf {Z} y _ {\alpha}.
\]

By Lemma 5 and Prop. 2, \(\mathcal{U}_{+},\mathcal{U}_{-},\mathcal{U}_{0}\) are orders in the \(\mathbf{Q}\) -algebras \(\mathrm{U}(\mathfrak{n}_{+}),\mathrm{U}(\mathfrak{n}_{-}),\mathrm{U}(\mathfrak{h})\) , respectively, and

\[
\binom{\pm h + q}{p} \in \mathcal {U} _ {0} \quad \text { for } h \in \mathcal {H}, q \in \mathbf {Z}, p \in \mathbf {N}.\tag{23}
\]

Put \(\mathcal{L} = \mathcal{U}_{-}. \mathcal{U}_{0}. \mathcal{U}_{+} \subset \mathrm{U}(\mathfrak{g})\) . By (22), \(\mathcal{L}\) is a lattice in \(\mathrm{U}(\mathfrak{g})\) . By construction,

\[
\mathcal {U} _ {-}. \mathcal {L} \subset \mathcal {L}\tag{24}
\]

\[
\mathcal {L}. \mathcal {U} _ {+} \subset \mathcal {L}\tag{25}
\]

while Lemma 3 and (23) imply that

\[
\mathcal {U} _ {0}. \mathcal {L} \subset \mathcal {L}\tag{26}
\]

\[
\mathcal {L}. \mathcal {U} _ {0} \subset \mathcal {L}.\tag{27}
\]

\[
\text {   Let   } \alpha \in \mathrm{B}, n \in \mathbf {N}, r \in \mathbf {N}, \varphi = (\varphi (i)) \in \mathrm{B} ^ {r}, \text {   and   }
\]

\[
(m (1), \dots , m (r)) \in \mathbf {N} ^ {r}.
\]

We show that

\[
x _ {\alpha} ^ {(n)} y _ {\varphi (1)} ^ {(m (1))} \dots y _ {\varphi (r)} ^ {(m (r))} \in \mathcal {L}\tag{28}
\]

or equivalently, in view of (25), that

\[
[ x _ {\alpha} ^ {(n)}, y _ {\varphi (1)} ^ {(m (1))} \dots y _ {\varphi (r)} ^ {(m (r))} ] \in \mathcal {L}.\tag{29}
\]

We argue by induction on \(r\) . The bracket to be studied is the sum of the terms

\[
y _ {\varphi (1)} ^ {(m (1))} \dots y _ {\varphi (k)} ^ {(m (k))} [ x _ {\alpha} ^ {(n)}, y _ {\varphi (k + 1)} ^ {(m (k + 1))} ] y _ {\varphi (k + 2)} ^ {(m (k + 2))} \dots y _ {\varphi (r)} ^ {(m (r))}.\tag{30}
\]

For \(\alpha \neq \varphi (k + 1)\) , \(x_{\alpha}\) and \(y_{\varphi (k + 1)}\) commute, so \([x_{\alpha}^{(n)},y_{\varphi (k + 1)}^{(m(k + 1))}] = 0\) . If \(\alpha = \varphi (k + 1)\) , the expression (30) is, by (17), the sum of expressions of the form

\[
y _ {\varphi (1)} ^ {(m (1))} \dots y _ {\varphi (k)} ^ {(m (k))} y _ {\varphi (k + 1)} ^ {(m (k + 1) - p)} \binom{q - h}{p} x _ {\alpha} ^ {(n - p)} y _ {\varphi (k + 2)} ^ {(m (k + 2))} \dots y _ {\varphi (r)} ^ {(m (r))}\tag{31}
\]

where \(q \in \mathbf{Z}, p \in \mathbf{N} - \{0\}, h \in \mathcal{H}\) . The induction hypothesis, together with (24) and (26), proves that the expression (31) belongs to \(\mathcal{L}\) . We have thus proved (28).

By (28), \(x_{\alpha}^{(n)}\mathcal{U}_{-}\subset \mathcal{L}\) ; thus, by (25) and (27), \(x_{\alpha}^{(n)}\mathcal{L}\subset \mathcal{L}\) , so \(\mathcal{U}_{+}.\mathcal{L}\subset \mathcal{L}\) and

\[
\mathcal {L}. \mathcal {L} \subset \mathcal {U} _ {-}. \mathcal {U} _ {0}. \mathcal {L} \subset \mathcal {U} _ {-}. \mathcal {L} \subset \mathcal {L}.
\]

Thus, \(\mathcal{L}\) is a \(\mathbf{Z}\) -subalgebra of \(\mathrm{U}(\mathfrak{g})\) , so \(\mathcal{U} = \mathcal{L}\) . If \(c\) is the coproduct of \(\mathrm{U}(\mathfrak{g})\) , \(c(\mathcal{U}) \subset \mathcal{U} \otimes_{\mathbf{Z}} \mathcal{U}\) (no. 2, Prop. 1). Let \(\gamma\) be the counit of \(\mathrm{U}(\mathfrak{g})\) . Since \(\gamma(x_{\alpha}^{(n)}) = \gamma(y_{\alpha}^{(n)}) = \gamma\left(\binom{h}{n}\right) = 0\) for \(n > 0\) , we have \(\gamma(\mathcal{U}) \subset \mathbf{Z}\) . This proves (i). On the other hand,

\[
\mathcal {U} \cap \mathfrak {h} = \mathcal {L} \cap \mathfrak {h} = \mathcal {U} _ {0} \cap \mathfrak {h} = \mathcal {H}
\]

by Prop. 2 of no. 4; similarly,

\[
\mathcal {U} \cap \mathfrak {g} ^ {\alpha} = \mathcal {U} _ {+} \cap \mathfrak {g} ^ {\alpha} = \mathbf {Z} x _ {\alpha}
\]

by Lemma 5. This proves (ii).

Remark 4. By Prop. 5 of §4, no. 4, there exists a unique automorphism \(\theta\) of \(\mathfrak{g}\) such that \(\theta(x_{\alpha}) = y_{\alpha}\) and \(\theta(y_{\alpha}) = x_{\alpha}\) for all \(\alpha \in \mathrm{B}\) , and \(\theta(h) = -h\) for all \(h \in \mathfrak{h}\) ; we have \(\theta^2 = 1\) . By construction of \(\mathcal{U}\) , we see that the automorphism of \(\mathrm{U}(\mathfrak{g})\) that extends \(\theta\) leaves \(\mathcal{U}\) stable.

COROLLARY 1. Put \(\mathcal{G} = \mathcal{U} \cap \mathfrak{g}\) . Then \(\mathcal{G}\) is an order in the Lie algebra \(\mathfrak{g}\) , stable under \(\theta\) . We have \(\mathcal{G} = \mathcal{H} + \sum_{\alpha \in \mathrm{R}} (\mathcal{G} \cap \mathfrak{g}^{\alpha})\) . For all \(\alpha \in \mathrm{B}\) and all \(n \in \mathbf{N}\) , the maps \((\operatorname{ad} x_{\alpha})^{n} / n!\) , \((\operatorname{ad} y_{\alpha})^{n} / n!\) leave \(\mathcal{U}\) and \(\mathcal{G}\) stable.

The first assertion is clear. The second follows by considering the graduation of type \(Q(R)\) on \(U(\mathfrak{g})\) and U. The third follows from Lemma 2 of no. 5.

COROLLARY 2. Let \(w \in \mathrm{W}\) . There exists an elementary automorphism \(\varphi\) of \(\mathfrak{g}\) that commutes with \(\theta\) , leaves \(\mathcal{G}\) and \(\mathcal{U}\) stable, and extends \(w\) .

It suffices to treat the case in which \(w\) is of the form \(s_{\alpha}\) ( \(\alpha \in \mathrm{B}\) ). Note first of all that \(\operatorname{ad} x_{\alpha}\) and \(\operatorname{ad} y_{\alpha}\) are locally nilpotent on \(\mathrm{U}(\mathfrak{g})\) , in other words that for all \(u \in \mathrm{U}(\mathfrak{g})\) there exists an integer \(n\) such that \((\operatorname{ad} x_{\alpha})^n u = (\operatorname{ad} y_{\alpha})^n u = 0\) . This enables us to define the automorphisms \(e^{\operatorname{ad} x_{\alpha}} = \sum_{n=0}^{\infty} \frac{1}{n!} (\operatorname{ad} x_{\alpha})^n\) and \(e^{\mathrm{ad}y_{\alpha}}\) of \(\mathrm{U}(\mathfrak{g})\) ; we verify immediately that these automorphisms of \(\mathrm{U}(\mathfrak{g})\) leave \(\mathcal{U}\) stable. Put \(\varphi_1 = e^{\mathrm{ad}x_\alpha}e^{\mathrm{ad}y_\alpha}e^{\mathrm{ad}x_\alpha}\) , \(\varphi_2 = e^{\mathrm{ad}y_\alpha}e^{\mathrm{ad}x_\alpha}e^{\mathrm{ad}y_\alpha}\) . We have \(\varphi_1|\mathfrak{g} = \varphi_2|\mathfrak{g}\) (§2, no. 2, formula (1)), so \(\varphi_1 = \varphi_2\) . Put \(\varphi_1 = \varphi_2 = \varphi\) . We have \(\theta \varphi \theta^{-1} = \varphi\) , so \(\theta\) and \(\varphi\) commute. On the other hand, \(\varphi |\mathfrak{h} = w\) by §2, no. 2, Lemma 1.

COROLLARY 3. Let \(\alpha \in \mathrm{R}\) . If \(x \in \mathcal{G} \cap \mathfrak{g}^{\alpha}\) and \(n \in \mathbf{N}\) , we have \(x^{(n)} \in \mathcal{U}\) , and \((\operatorname{ad} x)^n / n!\) leaves \(\mathcal{G}\) and \(\mathcal{U}\) stable.

This is clear if \(\alpha \in \mathrm{B}\) , by construction of \(\mathcal{U}\) and Cor. 1. In the general case, there exists \(w \in \mathrm{W}\) such that \(w(\alpha) \in \mathrm{B}\) (Chap. VI, §1, no. 5, Prop. 15). By Cor. 2, there exists an automorphism \(\varphi\) of \(\mathfrak{g}\) that leaves \(\mathcal{G}\) and \(\mathcal{U}\) stable and takes \(\mathfrak{g}^{\alpha}\) to \(\mathfrak{g}^{w(\alpha)}\) , hence the corollary by transport of structure.

COROLLARY 4. There exists a Chevalley system \((X_{\alpha})_{\alpha \in \mathrm{R}}\) in \((\mathfrak{g},\mathfrak{h})\) (§2, no. 4, Def. 3) such that \(X_{\alpha} = x_{\alpha}\) and \(X_{-\alpha} = y_{\alpha}\) for \(\alpha \in \mathrm{B}\) . For every Chevalley system \((X_{\alpha}^{\prime})_{\alpha \in \mathrm{R}}\) having these properties, and for all \(\alpha \in \mathrm{R}\) , \(X_{\alpha}^{\prime}\) is a basis of \(\mathcal{G} \cap \mathfrak{g}^{\alpha}\) .

For \(\alpha \in \mathrm{B}\) , put \(X_{\alpha} = x_{\alpha}, X_{-\alpha} = y_{\alpha}\) . For \(\alpha \in \mathrm{R}_{+} - \mathrm{B}\) , choose a \(w \in \mathrm{W}\) such that \(w(\alpha) \in \mathrm{B}\) and an automorphism \(\varphi\) of \(\mathfrak{g}\) such that \(\theta \varphi = \varphi \theta\) , \(\varphi(\mathcal{G}) = \mathcal{G}\) and \(\varphi(h) = w^{-1}(h)\) for \(h \in \mathfrak{h}\) (Cor. 2); put \(X_{\alpha} = \varphi(x_{w(\alpha)})\) , \(X_{-\alpha} = \varphi(y_{w(\alpha)})\) . Then

\[
[ X _ {- \alpha}, X _ {\alpha} ] = \varphi ([ y _ {w (\alpha)}, x _ {w (\alpha)} ]) = \varphi (H _ {w (\alpha)}) = w ^ {- 1} (H _ {w (\alpha)}) = H _ {\alpha}
\]

\[
\theta (X _ {\alpha}) = \theta \varphi (x _ {w (\alpha)}) = \varphi \theta (x _ {w (\alpha)}) = \varphi (y _ {w (\alpha)}) = X _ {- \alpha}
\]

so \((X_{\alpha})_{\alpha \in \mathrm{R}}\) is a Chevalley system. Moreover,

\[
\mathcal {G} \cap \mathfrak {g} ^ {\alpha} = \varphi (\mathcal {G} \cap \mathfrak {g} ^ {w (\alpha)}) = \varphi (\mathbf {Z} x _ {w (\alpha)}) = \mathbf {Z} X _ {\alpha}\tag{32}
\]

\[
\mathcal {G} \cap \mathfrak {g} ^ {- \alpha} = \varphi (\mathcal {G} \cap \mathfrak {g} ^ {- w (\alpha)}) = \varphi (\mathbf {Z} y _ {w (\alpha)}) = \mathbf {Z} X _ {- \alpha}.\tag{33}
\]

Let \((X'_{\alpha})_{\alpha \in \mathrm{R}}\) be a Chevalley system such that \(X'_{\alpha} = x_{\alpha}, X'_{-\alpha} = y_{\alpha}\) for \(\alpha \in \mathrm{B}\) . Let \(\mathrm{S}\) be the set of \(\alpha \in \mathrm{R}\) such that \(X'_{\alpha} = \pm X_{\alpha}\) . By §2, no. 4, Prop. 7, \(\mathrm{S}\) is a closed set of roots. Since \(\mathrm{S} \supset \mathrm{B} \cup (-\mathrm{B})\) , we have \(\mathrm{S} = \mathrm{R}\) (Chap. VI, §1, no. 6, Prop. 19). Thus, by (32) and (33), we have \(\mathcal{G} \cap \mathfrak{g}^{\alpha} = \mathbf{Z}X'_{\alpha}\) for all \(\alpha \in \mathrm{R}\) .

Remarks. 5) Let \((X_{\alpha})_{\alpha \in \mathrm{R}}\) be the Chevalley system constructed above. If \(\alpha, \beta, \alpha + \beta \in \mathrm{R}\) and if we put \([X_{\alpha}, X_{\beta}] = \mathrm{N}_{\alpha, \beta} X_{\alpha + \beta}\) , we have \([X_{\alpha}, X_{\beta}] \in \mathcal{G} \cap \mathfrak{g}^{\alpha + \beta}\) , and we recover the fact that \(\mathrm{N}_{\alpha, \beta} \in \mathbf{Z}\) (cf.~§2, no. 4, Prop. 7).

\begin{enumerate}
\def\labelenumi{\arabic{enumi})}
\setcounter{enumi}{5}
\tightlist
\item
  We have obtained in passing a new proof of the existence of Chevalley systems (cf.~§4, no. 4, Cor. of Prop. 5), independent of Lemma 4, §2.
\end{enumerate}

